% Chapter 1, Topic from _Linear Algebra_ Jim Hefferon
%  http://joshua.smcvt.edu/linearalgebra
%  2001-Jun-09
\topic{Keakuratan Perhitungan}
\index{keakuratan!Metode Gauss|(}
\index{Metode Gauss!keakuratan|(}
Metode Gauss mudah diimplementasikan pada komputer.
Kode di bawah ini memberikan ilustrasinya.
Kode tersebut bekerja pada matriks $\nbyn{n}$ bernama~\lstinline[style=inline]!a!, 
dengan melakukan kombinasi baris menggunakan baris pertama, kemudian
baris kedua, dan seterusnya.
\begin{lstlisting}
for(row=1; row<=n-1; row++){
  for(row_below=row+1; row_below<=n; row_below++){
    multiplier=a[row_below,row]/a[row,row];
    for(col=row; col<=n; col++){
      a[row_below,col]-=multiplier*a[row,col];
    }
  }
}
\end{lstlisting}
Kode ini ditulis dalam bahasa~C.\index{bahasa C}
Perulangan 
\lstinline[style=inline]!for(row=1; row<=n-1; row++){ .. }!
memberi nilai awal~$1$ pada \lstinline[style=inline]!row!, lalu berulang selama
\lstinline[style=inline]!row! kurang dari atau sama dengan $n-1$; pada setiap
putaran, nilai \lstinline[style=inline]!row! dinaikkan satu dengan operasi 
\lstinline[style=inline]!++!.
Konstruksi bahasa lain yang tidak langsung terlihat jelas adalah bahwa
\lstinline[style=inline]!-=! pada perulangan terdalam menghasilkan kode berikut.
\begin{lstlisting}[breaklines=true]
a[row_below,col]=-1*multiplier*a[row,col]+a[row_below,col]
\end{lstlisting}

Meskipun kode itu merupakan upaya awal untuk mengotomatisasi Metode Gauss,
kode tersebut masih naif.
Salah satu masalahnya adalah bahwa kode itu mengasumsikan entri pada posisi
\lstinline[style=inline]!row,row! tidak nol.
Jadi, salah satu perluasannya harus menangani kasus ketika ditemukannya nol
pada posisi tersebut mengharuskan pertukaran baris atau menghasilkan
kesimpulan bahwa matriks itu singular.

Kita dapat menambahkan beberapa pernyataan \lstinline[style=inline]!if! untuk
menangani kasus-kasus tersebut, tetapi kita akan membahas segi lain yang
membuat kode ini naif.
Kode ini rentan terhadap masalah yang timbul akibat ketergantungan komputer
pada aritmetika titik mengambang.

Sebagai contoh, di atas kita telah melihat bahwa sistem singular harus
ditangani sebagai kasus tersendiri.
Namun, sistem yang hampir singular juga perlu ditangani dengan hati-hati.
Perhatikan sistem berikut (digit tambahannya berada di tempat signifikan
kesembilan).
\begin{equation*}
   \begin{linsys}{2}
                   x &+ &2y &= &3\hfill\hbox{}  \\
     1.000\,000\,01x &+ &2y &= &3.000\,000\,01
   \end{linsys}
   \tag{$*$}
\end{equation*}
Dengan mengamatinya saja, kita mudah menemukan solusi $x=1$, $y=1$.
Komputer lebih sulit menanganinya.
Jika komputer merepresentasikan bilangan real hingga delapan tempat signifikan,
yang disebut \definend{presisi tunggal}\index{presisi tunggal}, 
maka secara internal komputer itu akan merepresentasikan persamaan kedua
sebagai $1.000\,000\,0x+2y=3.000\,000\,0$, sehingga digit di tempat
kesembilan hilang.
Alih-alih melaporkan solusi yang benar, komputer ini akan menganggap kedua
persamaan sama dan melaporkan bahwa sistem tersebut singular.

Untuk memperoleh gambaran intuitif tentang bagaimana komputer dapat memberikan
hasil yang menyimpang sejauh itu, perhatikan grafik sistem berikut.
\begin{center}
  \includegraphics{gr/mp/ch1.33}
\end{center}
Kita tidak dapat membedakan kedua garis tersebut;
sistem ini hampir singular dalam arti bahwa kedua garisnya hampir sama.
Hal ini membuat sistem~($*$) memiliki sifat bahwa perubahan kecil pada sebuah
persamaan dapat menyebabkan perubahan besar pada solusinya.
Misalnya, mengubah $3.000\,000\,01$ menjadi $3.000\,000\,03$ 
mengubah titik potong dari $(1,1)$ menjadi $(3,0)$.
Solusinya berubah drastis bergantung pada digit kesembilan, dan ini menjelaskan
mengapa komputer dengan delapan tempat mengalami kesulitan.
Masalah yang sangat peka terhadap ketidakakuratan atau ketidakpastian pada nilai
masukan disebut \definend{berkondisi buruk}.\index{masalah berkondisi buruk}

Contoh di atas memperlihatkan satu alasan sebuah sistem dapat sulit diselesaikan
dengan komputer.
Kelebihan contoh tersebut adalah bahwa gambar dua garis yang hampir sama
memberikan gambaran yang mudah diingat tentang salah satu cara terjadinya
kesulitan numerik.
Sayangnya, gambaran ini tidak berguna ketika kita ingin menyelesaikan suatu
sistem yang besar.
Biasanya kita tidak memahami geometri suatu sistem besar yang sembarang.

Selain kecilnya sudut antara beberapa permukaan linear, ada alasan lain yang
dapat membuat hasil komputer tidak dapat diandalkan.
Sebagai contoh, perhatikan sistem berikut (dari \cite{Hamming}).
\begin{equation*}
  \begin{linsys}{2}
     0.001x  &+  &y  &=  &1  \\
          x  &-  &y  &=  &0
  \end{linsys}
\tag*{($**$)}\end{equation*}
Persamaan kedua memberikan $x=y$, sehingga $x=y=1/1.001$ dan dengan demikian
nilai kedua variabel sedikit lebih kecil dari $1$.
Komputer yang menggunakan dua digit merepresentasikan sistem ini secara
internal sebagai berikut (demi kejelasan, kita akan mengerjakan contoh ini
dengan aritmetika titik mengambang dua digit, tetapi contoh serupa dengan
delapan digit atau lebih mudah dibuat).
\begin{equation*}
  \begin{linsys}{2}
    (1.0\times 10^{-3})\cdot x  &+  &(1.0\times 10^{0})\cdot y  &=  &1.0\times 10^{0}  \\
    (1.0\times 10^{0})\cdot x   &-  &(1.0\times 10^{0})\cdot y  &=  &0.0\times 10^{0}
  \end{linsys}
\end{equation*}
Langkah reduksi baris $-1000\rho_1+\rho_2$ menghasilkan persamaan kedua
$-1001y=-1000$, yang dibulatkan komputer ini menjadi dua tempat sebagai
$(-1.0\times 10^{3})y=-1.0\times 10^{3}$.
Dari persamaan kedua, komputer menyimpulkan bahwa $y=1$, lalu dari persamaan
pertama menyimpulkan bahwa $x=0$.
Nilai~$y$ cukup dekat, tetapi nilai~$x$ buruk\Dash
rasio antara jawaban sebenarnya dan jawaban komputer adalah
tak berhingga.\index{keakuratan!galat pembulatan}
Singkatnya, penyebab lain dari keluaran yang tidak dapat diandalkan adalah
ketergantungan komputer pada aritmetika titik mengambang ketika kode
penyelesaian sistem menggunakan entri utama yang kecil. 

Pemrogram berpengalaman mungkin menanggapinya dengan menggunakan
\definend{presisi ganda},\index{presisi ganda} %
yang mempertahankan enam belas digit signifikan, atau mungkin menggunakan
ukuran yang lebih besar lagi.
Cara ini memang akan menyelesaikan banyak masalah.
Namun, presisi ganda membutuhkan lebih banyak memori, dan kita jelas dapat
memodifikasi contoh di atas agar menimbulkan masalah yang sama pada digit
ketujuh belas, sehingga presisi ganda bukanlah solusi mujarab.
Kita memerlukan strategi untuk meminimalkan masalah numerik, sekaligus panduan
tentang seberapa jauh solusi yang dilaporkan dapat dipercaya.

Perbaikan mendasar terhadap kode naif di atas adalah tidak menentukan faktor
untuk kombinasi baris hanya dengan mengambil entri pada posisi
\lstinline[style=inline]!row,row!, tetapi dengan memeriksa semua entri pada
kolom \lstinline[style=inline]!row! di bawah entri
\lstinline[style=inline]!row,row! dan mengambil entri yang kemungkinan besar
memberikan hasil andal karena tidak terlalu kecil.
Cara ini disebut \definend{pemivotan parsial}.%
\index{pemivotan parsial}\index{pemivotan!parsial} 

Sebagai contoh, untuk menyelesaikan sistem bermasalah~($**$) di atas, mula-mula
kita memeriksa kedua persamaan untuk mencari entri terbaik yang akan digunakan,
dan memilih~$1$ pada persamaan kedua karena lebih mungkin memberikan hasil yang
baik.
Langkah kombinasi $-.001\rho_2+\rho_1$ menghasilkan persamaan pertama
$1.001y=1$, yang akan direpresentasikan komputer sebagai
$(1.0\times 10^{0})y=1.0\times 10^{0}$, sehingga diperoleh kesimpulan bahwa
$y=1$ dan, setelah substitusi balik, bahwa $x=1$; keduanya mendekati nilai yang
benar.  
Kita dapat menyesuaikan kode di atas untuk melakukan hal ini.
\begin{lstlisting}
for(row=1; row<=n-1; row++){
/* cari entri terbesar pada kolom ini (di baris max) */
  max=row;
  for(row_below=row+1; row_below<=n; row_below++){
    if (abs(a[row_below,row]) > abs(a[max,row]))
      max = row_below;
  }
/* tukar baris agar entri terbesar itu berpindah ke atas */
  for(col=row; col<=n; col++){
    temp=a[row,col];
    a[row,col]=a[max,col];
    a[max,col]=temp;
  }
/* lanjutkan seperti sebelumnya */
  for(row_below=row+1; row_below<=n; row_below++){
    multiplier=a[row_below,row]/a[row,row];
     for(col=row; col<=n; col++){
       a[row_below,col]-=multiplier*a[row,col];
    }
  }
}
\end{lstlisting}

Analisis lengkap mengenai cara terbaik untuk mengimplementasikan Metode Gauss
berada di luar cakupan buku ini (lihat \cite{Wilkinson65}), tetapi metode yang
dianjurkan oleh kebanyakan ahli mula-mula mencari entri terbaik di antara para
kandidat, lalu menskalakannya menjadi bilangan yang lebih kecil kemungkinannya
untuk menimbulkan masalah.
Cara ini disebut
\definend{pemivotan parsial berskala}.\index{pemivotan parsial berskala}\index{pemivotan!parsial!berskala}

Selain memberikan hasil yang kemungkinan besar dapat diandalkan, kebanyakan
kode yang disusun dengan baik juga akan memberikan
\definend{bilangan kondisi}%
\index{bilangan kondisi}\index{matriks!bilangan kondisi}
yang menyatakan faktor pengali yang dapat memperbesar ketidakpastian pada
bilangan masukan hingga menjadi ketidakakuratan pada hasil yang dikembalikan
(lihat \cite{Rice}).

Pelajarannya adalah bahwa fakta bahwa Metode Gauss selalu berhasil secara
teoretis, dan bahwa kode komputer mengimplementasikan metode itu dengan benar,
tidak dengan sendirinya berarti bahwa jawabannya dapat diandalkan.
Dalam praktik, selalu gunakan paket perangkat lunak yang dikembangkan dengan
cermat oleh para ahli untuk menanggulangi hal-hal yang dapat bermasalah.

\begin{exercises}
  \item 
    Dengan menggunakan dua digit signifikan, jumlahkan $253$ dan $2/3$.
    \begin{answer}
      Notasi ilmiah memudahkan kita menyatakan pembatasan dua digit:
      $.25\times 10^{3}+.67\times 10^{0}=.25\times 10^{3}$.
      Perhatikan bahwa penambahan $2/3$ tidak berpengaruh pada jumlahnya.
    \end{answer}
  \item 
    Soal perpotongan garis ini berbeda dengan contoh yang dibahas di atas.
    \begin{center}
      \vcenteredhbox{\includegraphics{gr/mp/ch1.34}}
      \qquad
      $\displaystyle \begin{linsys}{2}
            x &+ &2y &= &3  \\
            3x &- &2y &= &1
      \end{linsys}$
    \end{center}
    Tunjukkan bahwa dalam sistem ini, perubahan kecil pada bilangan-bilangannya
    hanya akan menghasilkan perubahan kecil pada solusi dengan mengubah
    konstanta pada persamaan bawah menjadi $1.008$, lalu menyelesaikannya.
    Bandingkan hasilnya dengan solusi sistem yang tidak diubah.
    \begin{answer}
      Reduksi
      \begin{equation*}
        \grstep{-3\rho_1+\rho_2}
        \begin{linsys}{2}
          x  &+  &2y   &=  &3  \\
             &   &-8y  &=  &-7.992
        \end{linsys}
      \end{equation*}
      memberikan \( (x,y)=(1.002,0.999) \).
      Jadi, untuk sistem ini, perubahan kecil pada konstanta hanya menghasilkan
      perubahan kecil pada solusi.
    \end{answer}
  \item 
    Perhatikan sistem berikut (\cite{Rice}).
    \begin{equation*}
      \begin{linsys}{2}
        0.000\,3x  &+  &1.556y  &=  &1.559 \\
        0.345\,4x  &-  &2.346y  &=  &1.108
      \end{linsys}
    \end{equation*}
    \begin{exparts*}
      \partsitem Selesaikan sistem tersebut.
      \partsitem Selesaikan sistem tersebut dengan
         melakukan pembulatan hingga empat digit pada setiap langkah. 
    \end{exparts*}
    \begin{answer}
      \begin{exparts}
        \partsitem Solusi dengan keakuratan penuh adalah $x=10$ dan $y=1$.
        \partsitem Reduksi empat digit
          \begin{equation*}
            \grstep{-(.3454/.0003)\rho_1+\rho_2}
            \begin{amat}{2}
              .0003  &1.556  &1.569  \\
              0      &-1794   &-1805
            \end{amat}
          \end{equation*}
          memberikan kesimpulan bahwa~$y=1.006$, yang tidak terlalu buruk,
          dan bahwa $x=12.21$. 
          Tentu saja, nilai ini berbeda dua puluh persen dari jawaban yang benar.
      \end{exparts}
    \end{answer}
  \item 
    Pembulatan di dalam komputer sering memengaruhi hasil.
    Anggaplah mesin Anda memiliki delapan digit signifikan.
    \begin{exparts}
      \partsitem Tunjukkan bahwa mesin akan menghitung
         $(2/3)+((2/3)-(1/3))$ sebagai nilai yang tidak sama dengan
         $((2/3)+(2/3))-(1/3)$.
         Jadi, aritmetika komputer tidak bersifat asosiatif.
      \partsitem Bandingkan versi komputer dari $(1/3)x+y=0$
        dan $(2/3)x+2y=0$.
        Apakah dua kali persamaan pertama sama dengan persamaan kedua?
    \end{exparts}
    \begin{answer}
      \begin{exparts}
        \partsitem Untuk bentuk pertama, mula-mula, $(2/3)-(1/3)$ adalah
          $.666\,666\,67-.333\,333\,33=.333\,333\,34$,
          sehingga
          $(2/3)+((2/3)-(1/3))=.666\,666\,67+.333\,333\,34=1.000\,000\,0$.

          Untuk bentuk lainnya, mula-mula
          $((2/3)+(2/3))=.666\,666\,67+.666\,666\,67=1.333\,333\,3$,
          sehingga
          $((2/3)+(2/3))-(1/3)=1.333\,333\,3-.333\,333\,33=.999\,999\,97$.
        \partsitem Persamaan pertama adalah
          $.333\,333\,33\cdot x+1.000\,000\,0\cdot y=0$,
          sedangkan persamaan kedua adalah
          $.666\,666\,67\cdot x+2.000\,000\,0\cdot y=0$.
      \end{exparts}
    \end{answer}
  \item 
    Kondisi buruk tidak hanya bergantung pada matriks koefisien.
    Contoh dari \cite{Hamming} ini menunjukkan bahwa kondisi buruk dapat timbul
    akibat interaksi antara ruas kiri dan ruas kanan sistem.
    Misalkan $\varepsilon$ adalah bilangan real kecil.
    \begin{equation*}
      \begin{linsys}{3}
        3x  &+  &2y           &+  &z            &=  &6\hfill   \\
        2x  &+  &2\varepsilon y  &+  &2\varepsilon z  
                                          &=  &2+4\varepsilon\hfill \\
         x  &+  &2\varepsilon y  &-  &\varepsilon z   
                                          &=  &1+\varepsilon\hfill
      \end{linsys}
    \end{equation*}
    \begin{exparts}
      \partsitem Selesaikan sistem tersebut dengan tangan.
        Perhatikan bahwa faktor-faktor $\varepsilon$ hanya dapat dibagi habis
        karena bagian-bagian bulat pada ruas kanan saling meniadakan tepat,
        demikian pula pada ruas kiri. 
      \partsitem Selesaikan sistem tersebut dengan tangan, dengan membulatkan
        hingga dua digit dan menggunakan $\varepsilon=0.001$.
%      \partsitem If you have access to a computer package for solving
%        linear systems, see if you can find an $\varepsilon$ small enough
%        to get your package to give incorrect results.
    \end{exparts}
    \begin{answer}
      \begin{exparts}
        \partsitem Perhitungan berikut
          \begin{align*}
            &\grstep[-(1/3)\rho_1+\rho_3]{-(2/3)\rho_1+\rho_2}
            \begin{amat}{3}
              3  &2                   &1                   &6               \\
              0  &-(4/3)+2\varepsilon &-(2/3)+2\varepsilon &-2+4\varepsilon \\
              0  &-(2/3)+2\varepsilon &-(1/3)-\varepsilon  &-1+\varepsilon 
            \end{amat}                                                    \\
            &\grstep{-(1/2)\rho_2+\rho_3}
            \begin{amat}{3}
              3  &2                   &1                   &6               \\
              0  &-(4/3)+2\varepsilon &-(2/3)+2\varepsilon &-2+4\varepsilon \\
              0  &\varepsilon         &-2\varepsilon       &-\varepsilon 
            \end{amat}
          \end{align*}
          menghasilkan persamaan ketiga $y-2z=-1$.
          Substitusi ke persamaan kedua memberikan
          $((-10/3)+6\varepsilon)\cdot z=(-10/3)+6\varepsilon$,
          sehingga $z=1$ dan kemudian $y=1$.
          Dengan nilai-nilai tersebut, persamaan pertama menyatakan bahwa
          $x=1$. 
        \partsitem 
          Seperti di atas, notasi ilmiah memudahkan kita menyatakan pembatasan
          jumlah digit.

          Solusi dengan mempertahankan dua digit adalah
          $z=2.1$, $y=.46$, dan $x=.99$.
          \begin{multline*}
            \begin{amat}{3}
              .30\times 10^{1}  &.20\times 10^{1}  &.10\times 10^{1} 
                 &.60\times 10^{1}        \\
              .20\times 10^{1}  &.20\times 10^{-3} &.20\times 10^{-3} 
                 &.20\times 10^{1}        \\
              .10\times 10^{1}  &.20\times 10^{-3}  &-.10\times 10^{-3} 
                 &.10\times 10^{1}        
            \end{amat}                                           \\
            \begin{aligned}
            &\grstep[-(1/3)\rho_1+\rho_3]{-(2/3)\rho_1+\rho_2}
            \begin{amat}{3}
              .30\times 10^{1}  &.20\times 10^{1}  &.10\times 10^{1} 
                 &.60\times 10^{1}        \\
              0                 &-.13\times 10^{1} &-.67\times 10^{0} 
                 &-.20\times 10^{1}        \\
              0                 &-.67\times 10^{0}  &-.33\times 10^{0} 
                 &-.10\times 10^{1}        
            \end{amat}                                          \\
            &\grstep{-(.67/1.3)\rho_2+\rho_3}
            \begin{amat}{3}
              .30\times 10^{1}  &.20\times 10^{1}  &.10\times 10^{1} 
                 &.60\times 10^{1}        \\
              0                 &-.13\times 10^{1} &-.67\times 10^{0} 
                 &-.20\times 10^{1}        \\
              0                 &0                  &.15\times 10^{-1} 
                 &.31\times 10^{-1}        
            \end{amat}
            \end{aligned}
          \end{multline*}
      \end{exparts}
    \end{answer}
\end{exercises}
\index{keakuratan!Metode Gauss|)}
\index{Metode Gauss!keakuratan|)}

\endinput

